\documentclass{article}
\usepackage{graphicx}
\usepackage[margin=1.5cm]{geometry}
\usepackage{amsmath}

\begin{document}
\twocolumn

\title{Chapter 4 Warm-Up: Boolean Logic and Decisions}
\author{Prof. Jordan C. Hanson}

\maketitle

\section{Boolean Variables}

\begin{enumerate}

\item A Boolean variable can have one of two values:
\texttt{True} or \texttt{False}.  Comparisons also produce
Boolean values.

\begin{itemize}

\item Enter the following statements in the interpreter:

\begin{verbatim}
is_raining = False
temperature = 72

print(is_raining)
print(temperature > 70)
print(temperature == 72)
print(temperature < 60)
\end{verbatim}

\end{itemize}
\end{enumerate}


\section{Boolean Operations and Decisions}

\begin{enumerate}

\item Python combines Boolean conditions using
\texttt{and}, \texttt{or}, and \texttt{not}.

\begin{itemize}

\item Predict the values of

\begin{verbatim}
print(True and True)
print(True and False)
print(True or False)
print(False or False)
print(not True)
\end{verbatim}

\item Now write a short program that collects information
about a student:

\begin{verbatim}
score = int(input("Score: "))

submitted = input(
    "Submitted? yes/no: ") == "yes"

attendance = input(
    "Attendance OK? yes/no: ") == "yes"
\end{verbatim}

Create the Boolean variable

\begin{verbatim}
ready = submitted and attendance
\end{verbatim}

and use it in the following chained decision:

\begin{verbatim}
if score >= 90 and ready:
    print("Excellent")
elif score >= 70 and ready:
    print("Pass")
elif not submitted or not attendance:
    print("Incomplete")
else:
    print("Revise")
\end{verbatim}

Run the program several times.  Find inputs that cause
each of the four branches to execute.

\end{itemize}
\end{enumerate}

\section{Truth-Table Exercise}

\begin{enumerate}

\item Associativity changes the grouping of operands, while
commutativity changes their order.  We will test whether these
changes affect a Boolean expression.  Consider the four expressions

\begin{verbatim}
E1 = (p and q) and (not r)
E2 = p and (q and (not r))
E3 = (not r) and (q and p)
E4 = p and q and not r
\end{verbatim}

Here, \texttt{E2} changes the association of the
\texttt{and} operations, \texttt{E3} changes both grouping
and ordering, and \texttt{E4} relies on Python's precedence
rules.

\begin{itemize}

\item There are eight possible combinations of three Boolean
variables.  Complete the truth table:

\begin{center}
\setlength{\tabcolsep}{3pt}
\begin{tabular}{ccc|cccc}
$p$ & $q$ & $r$ &
E1 & E2 & E3 & E4 \\ \hline
T & T & T & & & & \\
T & T & F & & & & \\
T & F & T & & & & \\
T & F & F & & & & \\
F & T & T & & & & \\
F & T & F & & & & \\
F & F & T & & & & \\
F & F & F & & & &
\end{tabular}
\end{center}

\item Check the table using Python.  For each row, assign
values to \texttt{p}, \texttt{q}, and \texttt{r}, calculate
all four expressions, and print the result.

\item Compare the four result columns.  Does changing the
\textbf{association} of the \texttt{and} operations change
the truth value?

\item Does changing the \textbf{order} of the Boolean
variables change the truth value?

\item Based on your truth table, what evidence do you have
that Boolean \texttt{and} is associative and commutative?

\end{itemize}
\end{enumerate}

\end{document}